Tout le chapitre repose sur une seule opération : la division euclidienne. La parité, les
critères de divisibilité, le pgcd, la décomposition en facteurs premiers, l'écriture
décimale — chacune de ces notions n'est qu'une lecture du reste dans une division bien
choisie.

C'est pourquoi les définitions sont posées d'abord, et pourquoi elles sont si courtes : un
nombre pair est un nombre de la forme $2k$, le chiffre des unités est le reste modulo dix,
la somme des chiffres se définit par récurrence en retirant un chiffre à la fois. Tout ce
qui suit s'en déduit par calcul.

Deux résultats méritent qu'on s'y arrête. Le \emph{crible d'Ératosthène} n'est pas ici une
méthode décrite mais un calcul effectué : la liste des nombres premiers inférieurs à cent
est produite par la machine, et l'égalité avec la liste attendue est vérifiée. Et
l'\emph{unicité de la décomposition en facteurs premiers}, que le programme admet, est
démontrée — au prix du lemme d'Euclide, qui n'est pas au programme du collège. L'écart est
signalé sur place.

Le chapitre se termine sur le calcul littéral, où la démonstration change de nature :
prouver qu'une identité vaut pour tout $x$ n'a rien à voir avec la vérifier sur quelques
valeurs. C'est la première fois de la scolarité que cette différence a des conséquences.
